\chapter{Vertical Structure and Vertical-Velocity Interpolation}\label{ch:vertical}

In a three-dimensional run the particle also moves vertically, driven by the model vertical velocity
$\Wvel$. This chapter describes the sign convention of $\Wvel$ and how each reader establishes it,
how the depth coordinate is organized for the two vertical-coordinate families, how the layer
containing a particle is found, how the horizontal velocity is read in depth --- per layer or linearly
between its data points --- how $\Wvel$ is interpolated to the particle in the vertical and then
in the horizontal, how a step advances the depth on a terrain-following mesh, and how a particle is
held within the water column. The two families are the fixed depth levels (z-levels) of FESOM2, FESOM~1.4
and the HEALPix products, and the terrain-following sigma coordinate of FESOM-C; they are implemented by
\emph{separate} routines throughout, and we give each explicitly. Depth $\depx$ and vertical velocity
$\Wvel$ are positive downward (\cref{sec:intro-conv}). On a z-level mesh the vertical position update
is
\begin{equation}\label{eq:vert-update}
  \depx_{n+1} = \depx_n + \Wvel\,\dt ,
\end{equation}
which deepens a particle wherever the flow is downwelling; on a sigma mesh the step advances the
particle's fraction of the water column instead (\cref{sec:vert-sigma-step}).

\section{The sign of the vertical velocity}\label{sec:vert-wsign}

Inside the engine the vertical velocity is the positive-down rate of the tracked depth. Ocean models
write the opposite sign, and every reader converts once, when a snapshot is read
(\cref{tab:vert-wsign}). The convention of a file is resolved in this order: the configuration key
\conf{mode.w\_positive} when it is \conf{up} or \conf{down}; otherwise the variable's CF
\code{positive} attribute; otherwise a CF \code{standard\_name} beginning with \code{upward} or
\code{downward}; otherwise the default of the reader's producer. The default, \conf{auto}, therefore
needs no setting for model output. For a FESOM2 dataset without a vertical velocity, the tool
\code{utils/reconstruct\_w.py} reconstructs one offline from $u$, $v$ and the sea-surface height, by
integrating the horizontal divergence down the column; it writes the result positive down, with the
attribute that says so.

\begin{table}[t]
\centering\footnotesize\setlength{\tabcolsep}{4pt}
\begin{tabular}{@{}p{30mm}p{20mm}p{62mm}p{16mm}@{}}
\toprule
producer (reader) & variable & file convention & conversion \\
\midrule
FESOM2 (\code{fesom2}) & \code{w} & positive up; in the zstar coordinate the velocity \emph{through} the moving interfaces & negated \\
HEALPix products (\code{healpix}) & \code{wo}, \code{avg\_wo} & positive up (``upward sea water velocity'') & negated \\
FESOM-C (\code{fesomc}) & \code{w} & positive up; the flux through the sigma interfaces per unit area, zero at the surface and the bottom & negated \\
offline reconstruction & \code{w} & positive down, attribute \code{positive = "down"} & none \\
\bottomrule
\end{tabular}
\caption{Vertical-velocity conventions of the supported producers and the conversion each reader
applies. The FESOM-C diagnostic \code{w\_cart} is refused: its flux term lacks the dia-sigma part of the
Cartesian vertical velocity.}
\label{tab:vert-wsign}
\end{table}

\codenote{Engine versions before September 2026 applied no conversion, so their three-dimensional
FESOM2 and HEALPix runs moved particles vertically in the wrong direction. The sign is now tested
against the equatorial upwelling in real model output. A vertical-velocity file written by the offline
reconstruction before that date carries no attribute and is read with \conf{w\_positive: down}.}

\paragraph{The zstar coordinate of FESOM2.} FESOM2's interfaces move with the sea surface, and its
\code{w} is the flow through them. The engine integrates it on the fixed reference levels, where the
horizontal velocities also live, which drops the factor $H/(H+\eta)$ of the zstar equation: a relative
error $\eta/H$ in the vertical rate, about $10^{-4}$ in the deep ocean and \SIrange{1}{2}{\percent} on a
shelf \SIrange{50}{100}{\metre} deep. The recorded depth is the zstar (reference) depth, which differs
from the physical depth by $\eta\,(1-\depx/H)$, at most $|\eta|$.

\codeanchor{\fn{resolve\_w\_sign} in \srcfile{readers/\_\_init\_\_.py}, applied by the FESOM2,
HEALPix and FESOM-C readers; the key \conf{mode.w\_positive} is validated in \srcfile{core/config.py}.}

\section{Vertical coordinate: z-levels and sigma}\label{sec:vert-coord}

The two families are distinguished only by how an interface depth is obtained, but the difference
propagates through every routine below. For the \emph{z-level} family the interface depths are a
fixed array $z_0=0<z_1<\dots<z_{K-1}$, the same at every horizontal location; the per-node sea floor is
the separate array \code{depth}. FESOM2 meshes hold the levels in the mesh file. FESOM~1.4 meshes take
them from \code{nod3d.out} or from the depth coordinate of level-wise data, and their per-node floor is
the depth of each column's last wet level (\cref{sec:mesh-fesom14}). A HEALPix product
takes them from \conf{mesh.levels} --- the path of a level table or an inline list of depths, required to
be strictly increasing --- or from the table a regional extract embeds; its per-node floor comes from
a bathymetry variable when the store has one, otherwise from the deepest level with finite data. For
the terrain-following \emph{sigma} family of FESOM-C the interface depths are reconstructed per water
column from a fixed set of fractions $\sigma_k$ and the local column thickness $H$,
\begin{equation}\label{eq:sigma}
  H = \max(d+\eta,\,D_{\min}), \qquad z_k = (1-\sigma_k)\,H,
\end{equation}
where $d$ is the still-water depth, $\eta$ the sea-surface height, and $D_{\min}$ a configurable minimum
thickness (default \SI{0.5}{\metre}). With this convention $\sigma=1$ maps to $z=0$ (the surface) and
$\sigma=0$ maps to $z=H$ (the floor). The key \conf{mode.max\_levels} keeps the top
\conf{max\_levels} interfaces of either family.

\codenote{The FESOM2 NetCDF and Zarr loaders store the z-levels positive down; the
\code{mesh.z\_levels} property takes the absolute value once more. Three-dimensional runs on FESOM2's
legacy ASCII mesh format are unsupported: its level reader takes the whole \code{aux3d.out} file, node depths
included, as the level table. (The FESOM~1.4 reader reads its own, different \code{aux3d.out} format
correctly, \cref{sec:mesh-fesom14}.) For sigma meshes the orientation $\sigma=1$ (surface) to $\sigma=0$ (floor) is
assumed by \cref{eq:sigma} but is \emph{not} validated when the mesh is read: a file written with the
opposite ordering would silently invert the column.}

On z-level meshes $\eta$ does not enter the vertical path at all: the layer search, the
vertical-velocity profile and the depth clamp depend only on the fixed levels and the per-node floor.
On sigma meshes $\eta$ sets the column thickness everywhere, and the column that matters is the one
\emph{at the particle}: every routine below uses $H_p$, the column thickness interpolated to the
particle position by the same horizontal method as the vertical velocity (\cref{sec:vert-whoriz}).

\section{Finding the layer}\label{sec:vert-layer}

Given a particle depth $\depx$, the engine identifies the layer that contains it, and the two families
do this differently. On a z-level mesh the search is over the fixed interfaces: a particle at or above
$z_0$ is assigned to the top layer, and otherwise the layer is the index $k$ with
$z_k\le\depx<z_{k+1}$. The search is bounded below by the local floor through a specific rule --- the
deepest admissible layer is $\max(k^\ast-1,0)$, where $k^\ast$ is the first index whose lower interface
$z_{k+1}$ exceeds the per-node bottom depth --- so a particle is never placed in a layer that lies below
its local sea floor. On a sigma mesh the search is over the interfaces $z_k=(1-\sigma_k)H_p$ of the
particle's own column; its deepest layer is simply $K-2$, bounded by the bottom interface, with no
analogue of the z-level step-back. This asymmetry is a genuine difference in behaviour between the two
families at the deepest layer, not merely a difference in formulae.

\codeanchor{Layer location is \fn{find\_layer}/\fn{find\_layers} (\srcfile{numerics/vertical.py},
z-level) and \fn{find\_layer\_sigma}/\fn{find\_layers\_sigma} (\srcfile{numerics/vertical\_sigma.py},
sigma, with the particle column passed as \code{H\_particle}); the per-particle layer index is
assembled by \fn{\_compute\_layers} (\srcfile{numerics/integrators.py}).}

\section{Horizontal velocity in depth: layers or linear}\label{sec:vert-uvdepth}

The horizontal velocity of a three-dimensional run is stored at a few discrete depths per column,
and the engine offers two rules for reading it in between, selected by \conf{mode.uv\_vertical}.
Most producers write $u$ and $v$ once per layer, at the layer centre: the FESOM2 node fields
\code{unod}/\code{vnod} and element fields \code{u}/\code{v}, the HEALPix \code{uo}/\code{vo}, and the
FESOM-C \code{u}/\code{v} on its sigma layers, $K-1$ rows each. Under the default, \conf{layer}, the
velocity is \emph{piecewise constant in depth}: every member of the horizontal stencil of \cref{ch:hvel}
--- a vertex or an element --- reads the row of the particle's layer, found by the rule of
\cref{sec:vert-layer} in the nearest node's column (on a sigma mesh, in the particle's own column), so a
particle that crosses a level sees the velocity jump. Under \conf{linear} the same data are read
\emph{linearly in depth} between the layer centres, and every member is read in its own column. FESOM~1.4
writes $u$ and $v$ on the levels themselves, and its fields are always read linearly in depth. Inside the
engine the rule is a property of the field, its \emph{placement} \code{VelocityField.uv\_vertical}, with the
three values of \cref{tab:vert-uv-placement}. This section gives the rule, the treatment of a stencil
member whose sea floor lies above the particle, its form on sigma meshes and for FESOM~1.4, and the
combinations the engine refuses.

\begin{table}[t]
\centering\footnotesize\setlength{\tabcolsep}{4pt}
\begin{tabular}{@{}>{\raggedright\arraybackslash}p{27mm}>{\raggedright\arraybackslash}p{24mm}>{\raggedright\arraybackslash}p{46mm}>{\raggedright\arraybackslash}p{40mm}@{}}
\toprule
placement & rows of $u$, $v$ & data points $p_k$ (positive down) & set by \\
\midrule
\code{"layer"} & $K-1$ & none: one value per layer & every reader; the default \\
\code{"layer\_linear"} & $K-1$, node or element layout & the layer centres $\tfrac12(z_k+z_{k+1})$; on a sigma mesh $1-\tfrac12(\sigma_k+\sigma_{k+1})$, in the sigma fraction & \conf{mode.uv\_vertical: linear}; the tracking layer marks every field it reads \\
\code{"level"} & $K$, node layout & the levels $z_k$ & the FESOM~1.4 reader; the option is refused there \\
\bottomrule
\end{tabular}
\caption{The three placements of the horizontal velocity in depth. Only \code{"layer"} is piecewise
constant; the other two are read by the depth-linear rule of this section.}
\label{tab:vert-uv-placement}
\end{table}

\paragraph{The data-point table.} For a depth-linear field the engine builds, once per mesh, placement
and layout, a table of the data points $p_0\le p_1\le\dots\le p_{M-1}$ ($M=K$ for levels, $M=K-1$ for
layer centres) and, for every \emph{column} $c$ --- a node for node-layout fields, an element for
element-layout fields --- two numbers: the index $\ell_c$ of its \emph{last usable point}, the deepest
point it may be read at, and its \emph{bed} $d_c$, the physical depth of its sea floor. For a node $i$,
\begin{equation}\label{eq:vert-uv-last}
  \ell_i = \max\Bigl(0,\ \min\bigl(M-1,\ \max(\{k : l_k\le d_i\}\cup\{0\}),\ D_i\bigr)\Bigr),
\end{equation}
where $l_k=p_k$ for level data and $l_k=z_{k+1}$, the layer's \emph{lower} level, for layer data;
$d_i$ is the node's bed \code{mesh.depth}, and $D_i$ an optional declaration by the reader of the last
finite point of its own data column (\code{mesh.metadata["uv\_data\_last"]}). For layer data this is the
rule of the layer search (\cref{sec:vert-layer}): a \emph{partial bottom layer}, one whose lower level lies
below the bed, is never read. The HEALPix reader declares where each column's $u$ data end, and the
FESOM~1.4 reader declares each column's last wet level; $D_i=-1$ means ``the surface only'', and the
surface value must then be finite. A NaN bed matches no level, which gives $\ell_i=0$; an infinite bed
follows the formula as written ($+\infty$ keeps every point, $-\infty$ the surface point only). For an
element $e$ the node depths cannot tell where its water ends --- FESOM2 fixes the number of levels of each
element first and gives a node the deepest of its elements --- so the element table takes both numbers
from the reader: $\ell_e=\max(0,\min(M-1,D_e))$ with $D_e$ the element's last wet layer
(\code{uv\_data\_last\_elem}), and $d_e$ its bed (\code{elem\_bed}), which the FESOM2 reader derives from
the mesh diagnostic (\cref{sec:mesh-elem-bottoms}). A \conf{mode.max\_levels} cut caps every $\ell_c$ at
the last kept point and leaves the bed unchanged. Points below $\ell_c$ are \emph{never read}, not even
multiplied by zero, so dry or fill values stored there are harmless.

\codeanchor{The table is \fn{depth\_points} with the frozen dataclass \fn{DepthPoints} (fields \code{p},
\code{last}, \code{bed}) and the builder \fn{\_build\_depth\_points}, in \srcfile{numerics/vertical\_uv.py}.
It is cached on the mesh per placement and layout and rebuilt whenever \code{mesh.sigma}, the bed array or
the declaration is replaced; the engine builds it from its own read-only snapshot of the declaration
(\fn{\_snapshot\_declaration}), which replaces the reader's array in \code{mesh.metadata}.}

\paragraph{Bracket and blend.} At a column $c$ and particle depth $z$ the depth is first clipped into the
column's usable range, the bracketing points $p_k\le z_c<p_{k+1}$ (with $k+1\le\ell_c$) are found, and
the stored values $f$ ($u$ or $v$) are blended:
\begin{equation}\label{eq:vert-uv-blend}
  z_c = \max\bigl(p_0,\ \min(z,\ p_{\ell_c})\bigr),\qquad
  f_c(z) = (1-\interpfrac)\,f_{k,c} + \interpfrac\,f_{k+1,c},\qquad
  \interpfrac = \frac{z_c-p_k}{p_{k+1}-p_k}.
\end{equation}
Above the first point and below the last usable point the value of that point is used: the ends are
constant, and nothing is extrapolated. A depth exactly on a point, at either end, or in a degenerate
interval (a spacing below $10^{-12}$) reads the point directly, without touching its neighbour; of two
coincident points the deeper one is read, except at the surface end, where a depth at or above $p_0$
reads point~0. Single-precision data are widened to double precision before the blend. Because the
bracket depends only on the column's own $\ell_c$, a member whose last usable point lies above the
particle --- a partial bottom layer, a \conf{max\_levels} cut, a declared prefix that ends above the bed ---
contributes the value of that point, provided it holds water at $z$ (\cref{fig:vert-uv}).

\codeanchor{\fn{uv\_bracket} returns $(k,\interpfrac)$, with $\interpfrac<0$ meaning ``read point $k$
directly'', and \fn{uv\_at} applies \cref{eq:vert-uv-blend}; both are compiled kernels in
\srcfile{numerics/vertical\_uv.py}.}

\begin{figure}[t]
\centering
\begin{tikzpicture}[scale=1.15]
  % levels z_0..z_5
  \foreach \y in {0,-0.8,-1.6,-2.4,-3.2,-4.0} \draw[interface] (0.1,\y)--(5.75,\y);
  \node[font=\footnotesize,left] at (0.1,0) {$z_0=0$};
  \node[font=\footnotesize,left] at (0.1,-0.8) {$z_1$};
  % beds: a thick floor segment and hatching below each column
  \foreach \x/\b in {0.8/-4.0,3.4/-1.9,5.0/-2.95}{
    \fill[pattern=north east lines,pattern color=gray!45] (\x-0.42,\b) rectangle (\x+0.42,-4.35);
    \draw[laggreen,line width=1.0pt] (\x-0.42,\b)--(\x+0.42,\b);
    \draw[gray!60] (\x,0)--(\x,\b);
    \node[meshnode] at (\x,0) {};}
  \node[font=\footnotesize,align=center,anchor=south] at (0.8,0.08) {vertex $A$\\wet};
  \node[font=\footnotesize,align=center,anchor=south] at (3.4,0.08) {vertex $B$\\dry at $z$};
  \node[font=\footnotesize,align=center,anchor=south] at (5.0,0.08) {vertex $C$\\wet};
  % data points (layer centres) down to each column's last usable point
  \foreach \y in {-0.4,-1.2,-2.0,-2.8,-3.6} \node[uvpoint] at (0.8,\y) {};
  \foreach \y in {-0.4,-1.2} \node[uvpoint] at (3.4,\y) {};
  \foreach \y in {-0.4,-1.2,-2.0} \node[uvpoint] at (5.0,\y) {};
  \node[uvpointopen] at (5.0,-2.8) {};
  \node[font=\footnotesize,left] at (0.68,-0.4) {$p_0$};
  \node[font=\footnotesize,left] at (0.68,-2.0) {$p_k$};
  \node[font=\footnotesize,left] at (0.68,-2.8) {$p_{k+1}$};
  % the particle depth
  \draw[gray!60,densely dotted,line width=0.7pt] (0.1,-2.62)--(5.75,-2.62);
  \node[drifter] at (2.5,-2.62) {};
  \node[font=\footnotesize,kelp!70!black,anchor=north] at (2.5,-2.74) {depth $z$};
  % A: the bracket
  \draw[<->,gray!80] (1.1,-2.0)--(1.1,-2.62) node[midway,right,font=\scriptsize,text=black]{$\interpfrac\,(p_{k+1}-p_k)$};
  % B: a wall at this depth
  \node[kelp!80!black,font=\large] at (3.4,-2.62) {$\times$};
  \node[font=\scriptsize,fill=white,inner sep=1pt,align=center] at (3.4,-3.6) {bed above $z$:\\a wall};
  % C: clipped to its last usable point
  \draw[-{Stealth[length=4pt]},lagnode,line width=0.8pt] (5.0,-2.62)--(5.0,-2.1);
  \node[font=\scriptsize,anchor=west] at (5.12,-2.15) {reads $p_{\ell_C}$};
  \node[font=\scriptsize,anchor=west] at (5.12,-2.79) {partial layer, never read};
  % what each column contributes
  \node[font=\scriptsize,anchor=north] at (0.8,-4.42) {$(1-\interpfrac)f_{k}+\interpfrac f_{k+1}$};
  \node[font=\scriptsize,anchor=north] at (3.4,-4.42) {left out};
  \node[font=\scriptsize,anchor=north] at (5.0,-4.42) {$f_{\ell_C}$};
  % depth axis
  \draw[-{Stealth[length=5pt]},gray] (-0.85,0.05)--(-0.85,-4.2);
  \node[gray,rotate=90,anchor=south,font=\footnotesize] at (-1.0,-2.1){depth (positive down)};
\end{tikzpicture}
\caption{The depth-linear rule at one particle (star, depth $z$) whose containing triangle has the
vertices $A$, $B$, $C$, drawn as three columns side by side. Diamonds are the data points --- here layer
centres between the levels (grey lines) --- down to each column's last usable point $p_{\ell_c}$; green
segments are the beds. Vertex $A$ is wet at $z$ and blends its two bracketing points. The bed of $B$ lies
above $z$, so $B$ is dry there: a wall, left out of the sum. Vertex $C$ is wet, but its bottom layer is cut
by its bed (a partial layer, open diamond, never read), so it contributes the value of its last usable
point. The node-layout result is $U=\tdir(w_A^{+}u_A+w_C^{+}u_C)/(w_A^{+}+w_C^{+})$
(\cref{eq:vert-uv-bary}); for element layout the columns are the elements of the stencil, under the same
rule.}
\label{fig:vert-uv}
\end{figure}

\paragraph{Walls at depth.} A column is \emph{wet} at the particle depth when $z\le d_c$, with $d_c$ the
node's \code{mesh.depth} --- the bed the end-of-step clamp of \cref{sec:vert-clamp} uses --- or the
element's \code{elem\_bed}. A NaN particle depth is read as $+\infty$, the layer path's reading of a NaN
depth, and is therefore dry wherever the bed is finite. A member whose bed lies above the particle is
\emph{dry}, and it is treated as a wall, like the coast: it is left out of the interpolation. We emphasize
that this is neither a zero velocity, which would brake every particle travelling along the wall, nor the
dry member's bottom value read below its bed, which could accelerate it. Each interpolator applies the
rule the way it treats the coast (\cref{tab:vert-uv-walls}). The node-layout barycentric scheme leaves a
dry vertex out of the sum of \cref{eq:bary-blend} and renormalises the weights of the wet vertices,
\begin{equation}\label{eq:vert-uv-bary}
  U = \tdir\,\frac{\sum_{i\ \text{wet}} \bwt{i}^{+}\,u_i(z)}{\sum_{i\ \text{wet}} \bwt{i}^{+}},
  \qquad \bwt{i}^{+}=\max(\bwt{i},0),
\end{equation}
where a negative weight --- a particle just outside an edge, which the containment test accepts down to
$-10^{-6}$ --- is counted as zero, so the result stays a convex blend of the wet vertex values; $V$ is
formed alike. The least-squares scheme leaves a dry element out of the plane fit \cref{eq:leastsq} over
the nearest node's ring, as the elements beyond a coast are absent from a coastal node's ring. The edge
least-squares scheme makes the edge to a dry neighbour a boundary edge, whose ghost \cref{eq:edge-ghost}
is built from the own element's value at $z$, mirrored (free slip) or negated (no slip); with no slip a wall
at depth therefore brakes the flow as the coast does. The element-layout barycentric scheme leaves a dry
element out of each vertex's area-weighted mean, and a vertex without a wet element out of the sum, the
weights of the others renormalised as in \cref{eq:vert-uv-bary}.

When every member is wet, the result is exactly the layer interpolator's arithmetic --- metric, summation
order and singular guard --- applied to the depth-linear values $f_c(z)$; a fit that is singular with every
member wet, such as a coastal ring, still fails the particle as in layer mode. Fallbacks apply only when
at least one member is dry, and only in three situations. First, when the wet members cannot support the
interpolator, the least-squares scheme (fewer than three wet elements, or a singular fit over them --- one
or two wet elements in a narrow trench, where the fit would otherwise fail and freeze the particle) and
the edge least-squares scheme (a singular fit over the own element and its wet neighbours) take the
area-weighted mean of the wet members, $\sum_e A_e f_e(z)/\sum_e A_e$, with the element areas of
\cref{sec:hvel-quad}; the barycentric scheme, in either layout, takes the equal-weight mean of the
remaining vertex values when their weights sum below $10^{-12}$ (the particle on a left-out vertex).
Second, when no member is wet --- possible, because the end-of-step clamp uses the nearest node's bed
only --- the node-layout scheme reads the vertex with the largest bed (the first on ties) and the element
schemes read the containing element, each at its last usable point; the least-squares scheme, whose ring
does not need the containing element, searches for it in this case only --- a particle outside the mesh
fails, and a containing element that the search finds outside the ring and wet at $z$ is read at $z$. Third, when the containing element of the edge least-squares scheme is itself dry, the scheme
takes the area-weighted mean of its wet edge neighbours, without ghosts, and with none wet the containing
element at its last usable point.

\begin{table}[t]
\centering\footnotesize\setlength{\tabcolsep}{4pt}
\begin{tabular}{@{}>{\raggedright\arraybackslash}p{24mm}>{\raggedright\arraybackslash}p{40mm}>{\raggedright\arraybackslash}p{38mm}>{\raggedright\arraybackslash}p{34mm}@{}}
\toprule
interpolator & a dry member & wet members insufficient & no member wet \\
\midrule
\conf{barycentric}, node layout & vertex left out; wet weights renormalised, negative ones as zero & weights sum $<10^{-12}$: equal-weight mean of the wet vertices & the vertex with the largest bed, at $p_{\ell}$ \\
\conf{leastsq} & element left out of the ring fit & fewer than three wet, or a singular fit: area-weighted mean of the wet elements & the containing element (searched), at $p_{\ell}$ \\
\conf{edge\_leastsq} & its edge a boundary edge with the \conf{edge\_leastsq\_bc} ghost & a singular fit: area-weighted mean of the wet members; the own element dry: of its wet neighbours & the containing element, at $p_{\ell}$ \\
\conf{barycentric}, element layout & element left out of each vertex mean; a vertex without one left out & weights sum $<10^{-12}$: equal-weight mean & the containing element, at $p_{\ell}$ \\
\bottomrule
\end{tabular}
\caption{The wall rule and its fallbacks per interpolator. With every member wet none of the three
right-hand columns applies and the result is the layer interpolator's arithmetic on the depth-linear
values.}
\label{tab:vert-uv-walls}
\end{table}

The velocity therefore jumps where a particle crosses the depth of a neighbour's bed, because that member
leaves or enters the stencil --- the analogue of the jump at a layer interface in layer mode. Near the bed
the two rules differ by more than the vertical interpolation. Layer mode reads every member in the row of
the nearest node's layer, a vertex or an element shallower than that layer included; FESOM2 files store
zeros there, so a shallow neighbour pulls the velocity --- for the least-squares scheme, the fitted plane
--- towards zero. Linear mode leaves that member out. For instance, a uniform flow along a trench, below
the bed of the shelf beside it, keeps its speed in linear mode and is slowed in layer mode
(\cref{sec:val-depth-linear}).

\paragraph{Time and stages.} In a time-interpolated run every stored value is first blended in time
(linear or PCHIP, \cref{ch:temporal}), then read in depth by \cref{eq:vert-uv-blend}, and then combined
horizontally. PCHIP is not linear, so this order is part of the contract: the fused kernels follow it and
agree bit for bit with the unfused path, which materializes the time blend first. A dry member is never
read, not even blended in time. Every Runge--Kutta stage reads the horizontal velocity at its own stage
depth --- Euler at the step-entry depth, RK2 at the entry and the midpoint depth, RK4 at its four stage
depths --- before any clamp: a stage below a member's bed meets that wall, a stage below every bed of its
stencil takes the no-member-wet fallback, and a stage above the surface reads $p_0$ everywhere.

\codeanchor{\raggedright The node-layout core is \fn{\_interp\_barycentric\_core\_dl} with the wall rule
\fn{\_dl\_node\_wall}; the element-layout cores are \fn{\_interp\_leastsq\_core\_dl},
\fn{\_interp\_edge\_leastsq\_core\_dl} and \fn{\_interp\_bary\_elem\_core\_dl}, which share the fallback
\fn{\_dl\_elem\_wet\_mean}; the wet test is \fn{uv\_wet} (\srcfile{numerics/vertical\_uv.py}). The fused
twins are \fn{\_fused\_euler\_step\_3d\_dl} (element layout, least squares, Euler) and
\fn{\_fused\_rk2\_step\_3d\_bary\_tinterp\_dl} and \fn{\_fused\_rk2\_step\_3d\_bary\_pchip\_dl} (node
layout, RK2, linear and PCHIP in time), in \srcfile{numerics/integrators.py}, with the time-blend vertex
cores \fn{\_interp\_barycentric\_core\_tinterp\_dl} and \fn{\_interp\_barycentric\_core\_pchip4\_dl}.
Element-layout RK2 and RK4 run unfused, as in layer mode.}

\paragraph{Sigma meshes.} A sigma column has the same layers everywhere, stretched to its own height, so
on a FESOM-C mesh the coordinate of the rule is the particle's \emph{sigma fraction} rather than its
depth --- the fraction of \cref{eq:vert-sigma-frac}, written $s$ here as in the code --- and the data points
are the layer centres in that fraction:
\begin{equation}\label{eq:vert-uv-sigma}
  s = \min\bigl(\max(z/H_p,\,0),\,1\bigr),\qquad p_k = 1-\tfrac12(\sigma_k+\sigma_{k+1}),\quad k=0,\dots,K-2,
\end{equation}
with the interfaces at $s_k=1-\sigma_k$ (0 at the surface, 1 at the bed). Only \code{"layer\_linear"}
data are served, on either layout. Every layer of every column holds water, so the table reads no
bathymetry, no declaration and no element bottom: $\ell_c=M-1$ and $d_c=+\infty$ for every column.
Every member is therefore wet at every $s$; there are no walls at depth and none of their fallbacks, and
the record check below requires every layer of every column to be finite. Every member is read at the
particle's $s$, that is, along the sigma surface through the particle --- as the vertical velocity is
(\cref{fig:vert-sigma}) and as layer mode reads every member in the row of the particle's sigma layer.
Each Runge--Kutta stage takes its $s$ from its own position, time and column: the integrators hand the
interpolator $s=\min(\max(z_{\text{stage}}/H_p(\pos_{\text{stage}},t_{\text{stage}}),0),1)$, with $H_p$
interpolated at the stage position from the stage's surface elevation (Euler at the entry, RK2 at the
entry and the midpoint, RK4 at all four stages), and keep metres for the vertical velocity, its layers
and the depth clamp. A NaN fraction is read as $+\infty$, that is, the bottom centre, and a
\conf{max\_levels} cut keeps the first interfaces from the surface, so every column reads down to the last
kept centre. On a sigma mesh ``exact'' means exact for a field linear in $s$. On a flat bed that is also
linear in $z$; on a sloping bed, however, a field linear in $z$ is not claimed, since sigma data cannot
represent it exactly under either rule.

\codeanchor{The sigma branch of \fn{\_build\_depth\_points} (\srcfile{numerics/vertical\_uv.py}); the
stage fraction is \fn{\_sigma\_frac\_for\_uv} with the stage column \fn{\_sigma\_col\_H}, in the sigma
branches of the unfused steps of \srcfile{numerics/integrators.py}. The sigma path stays unfused, as in
layer mode (\cref{ch:fused}).}

\paragraph{FESOM~1.4: velocities on the levels.} FESOM~1.4 is a P1 finite-element model: $u$, $v$, $\Wvel$,
temperature and salinity share the same three-dimensional nodes, which sit on the levels, and vary
linearly over the tetrahedra of its three-dimensional mesh. Its reader returns $u$ and $v$ with $K$ rows,
on the levels, in node layout and with the placement \code{"level"}, so every three-dimensional run on
FESOM~1.4 output is depth-linear between the levels, down to each column's last wet level, which the
reader declares as $D_i$ in \cref{eq:vert-uv-last}. The option \conf{mode.uv\_vertical: linear} is refused
for it, with the advice to remove the option. The engine does not interpolate on the model's tetrahedra:
it combines the barycentric weights of the surface triangle with the linear profile in depth of each
vertex's own column. The two agree on the levels and along each node's column. Inside the prism between
two levels they differ: the model's field is linear within each tetrahedron, whereas the engine's is
barycentric at every depth and linear in depth at every position, so it holds products of the two. We
therefore call it a \emph{prism-wise approximation} of FESOM~1.4's own interpolation. The vertical
velocity of FESOM~1.4 is also stored on the levels, which serve as the interfaces of
\cref{sec:vert-wprofile}.

\paragraph{The record check.} Every record a run reads with a depth-linear field is checked once per field
and mesh table, before the first read: a non-finite $u$ or $v$ at a point $k\le\ell_c$ of any column stops
the run, and the message names the component, the node or element, the point and the record. A lazy time
blend (\cref{ch:temporal}) is checked through its two parents, before anything is blended. Points below
$\ell_c$ may hold anything. With the option the HEALPix reader checks the raw record before it zeroes
missing values, so a surface mask that drifts between records is refused rather than zeroed.

\paragraph{Refusals.} The option \conf{linear} is refused for two-dimensional tracking
(\conf{Fdrift3D: false}); for a mesh that is neither z-level nor sigma; for an interpolator outside the
layout's row --- node layout takes \conf{barycentric} only, element layout \conf{leastsq},
\conf{edge\_leastsq} or \conf{barycentric}; for the \conf{rk45} integrator, which is two-dimensional only
(the step's entry gate refuses a depth-linear field for it as well); for a reader whose $u$ and $v$ sit
on the levels (FESOM~1.4); for element-layout data on a z-level mesh without element bottoms (a mesh
diagnostic with neither \code{nlevels} nor \code{zbar\_e\_bottom}, an ASCII mesh directory, or a region
extracted before the extraction tool kept them, \cref{sec:mesh-elem-bottoms}); and for element-layout data
on a mesh with ice-shelf cavities (\code{ulevels}${}>1$ anywhere). Ice-shelf cavities are not supported in
the code yet and are planned. Every one of these refusals comes before the output file exists: the value of
the key and the two-dimensional case when the configuration is loaded; the mesh type, the interpolator and
the reader right after the mesh is loaded; the layout and the element bottoms at the metadata-only probe
that precedes the writer (\cref{sec:hvel-leastsq}).
A bad record is refused when it is read. A field already marked depth-linear that reaches a run with the
option off is refused too, and so are the fields of one step whose placements differ.

\codeanchor{\raggedright The key is validated in \srcfile{core/config.py}; \fn{\_check\_uv\_vertical\_option} (after the mesh load), \fn{\_validate\_uv\_compat} and
\fn{\_check\_linear\_element\_mesh} (the pre-writer probe), and \fn{\_mark\_depth\_linear}, which marks
every field the velocity cache stores, in \srcfile{tracking.py}; \fn{refuse\_elem\_cavities} and
\fn{check\_depth\_linear\_record} in \srcfile{numerics/vertical\_uv.py}; the step-entry gate
\fn{\_check\_depth\_linear\_entry} in \srcfile{numerics/integrators.py}.}

\paragraph{Exactness.} Where every member is wet the rule reproduces the layer answer where it should:
data uniform in depth give the layer result bit for bit wherever every member reads a data point directly,
and between points within a few units in the last place of the largest member value (2 for node-layout
barycentric, 4 for element-layout barycentric, 8 for the two least-squares schemes; absolute bounds,
since member values of opposite sign may cancel). A field linear in depth is reproduced exactly between the
first and the last centre of the columns read, and a smooth profile converges at second order in the
layer thickness, layer mode at first order (\cref{sec:val-depth-linear}).

\section{The vertical velocity profile}\label{sec:vert-wprofile}

Within the bracketing interfaces the vertical velocity is interpolated linearly in \emph{physical
depth} (\cref{fig:vert}), for both families: on a sigma mesh the bracketing interface depths are first
reconstructed from \cref{eq:sigma} and the interpolation fraction is then formed in metres. With
\begin{equation}\label{eq:vert-alpha}
  \interpfrac = \frac{\depx-z_k}{z_{k+1}-z_k}\in[0,1],
\end{equation}
the profile value at a single node is
\begin{equation}\label{eq:vert-wlin}
  \Wvel(\depx) = \tdir\bigl[(1-\interpfrac)\,\Wvel_k + \interpfrac\,\Wvel_{k+1}\bigr],
\end{equation}
where $\Wvel_k,\Wvel_{k+1}$ are the interface values and $\tdir$ is the integration-direction factor. A
particle at or above the surface interface takes the surface value and a particle at or below the
deepest interface of the column takes the deepest value; no vertical extrapolation is performed, and a
degenerate layer (thickness below $10^{-12}$\,m) returns the upper-interface value. The local floor
enters through the horizontal combination, which evaluates each neighbour no deeper than its own
bottom (\cref{sec:vert-whoriz}).

On z-level meshes the bracket $(k,\interpfrac)$ depends only on the depth, so the least-squares and
inverse-distance combinations compute it once for the particle and reuse it at every neighbour deep
enough to contain the particle's depth; only shallower neighbours are bracketed again at their own
bottom. The result is bit-identical to bracketing each neighbour separately and makes the vertical-velocity
interpolation 1.1 to 1.4 times faster, the more so the more levels a column has.

\begin{figure}[t]
\centering
\begin{tikzpicture}[scale=1.15]
  \def\xl{0}\def\xr{4.0}
  \foreach \y in {0,-0.8,-1.6,-2.7,-3.5} \draw[interface] (\xl,\y)--(\xr,\y);
  \node[font=\footnotesize,left] at (\xl,0)    {$z_0=0$};
  \node[font=\footnotesize,left] at (\xl,-1.6) {$z_k$};
  \node[font=\footnotesize,left] at (\xl,-2.7) {$z_{k+1}$};
  \node[interfacept] at (1.3,-1.6) {}; \draw[velarrow] (1.3,-1.6)--++(0,-0.34);
  \node[font=\footnotesize,right] at (1.42,-1.44) {$\Wvel_k$};
  \node[interfacept] at (1.3,-2.7) {}; \draw[velarrow] (1.3,-2.7)--++(0,-0.34);
  \node[font=\footnotesize,right] at (1.42,-2.54) {$\Wvel_{k+1}$};
  \node[drifter] (D) at (2.55,-2.05) {};
  \node[font=\footnotesize,kelp!70!black,below right=2pt and 0pt] at (2.55,-2.05) {drifter, depth $\depx$};
  \draw[<->,gray!80] (4.9,-1.6)--(4.9,-2.05) node[midway,right,font=\footnotesize]{$\interpfrac\,\Delta z$};
  \draw[<->,gray!80] (4.9,-2.05)--(4.9,-2.7) node[midway,right,font=\footnotesize]{$(1-\interpfrac)\Delta z$};
  \draw[gray!50,densely dotted] (2.55,-2.05)--(4.9,-2.05);
  \draw[dashed,laggreen,line width=0.9pt] (\xl,-3.15)--(\xr,-3.15);
  \node[font=\footnotesize,laggreen,right] at (\xr,-3.18) {local floor};
  \fill[pattern=north east lines,pattern color=gray!45] (\xl,-3.15) rectangle (\xr,-3.55);
  \draw[-{Stealth[length=5pt]},gray] (-1.45,0.05)--(-1.45,-3.4);
  \node[gray,rotate=90,anchor=south,font=\footnotesize] at (-1.65,-1.7){depth $\depx$ (positive down)};
\end{tikzpicture}
\caption{Vertical interpolation of the vertical velocity, identical in form for the z-level and sigma
families. The interface values $\Wvel_k,\Wvel_{k+1}$ (open circles; downward arrows mark the
positive-down sign) bracket the drifter at depth $\depx$. The fraction
$\interpfrac=(\depx-z_k)/(z_{k+1}-z_k)$ is formed in physical depth, and the profile value is
$(1-\interpfrac)\Wvel_k+\interpfrac\Wvel_{k+1}$. No extrapolation is performed beyond the ends of the
column.}
\label{fig:vert}
\end{figure}

\codeanchor{The single-node profiles are \fn{\_interp\_w\_linear\_single} (z-level; the reusable
bracket is \fn{\_w\_bracket} with \fn{\_w\_at\_bracket}) and \fn{\_interp\_w\_sigma\_linear\_single}
(sigma), in \srcfile{numerics/vertical.py} and \srcfile{numerics/vertical\_sigma.py}.}

\section{Horizontal combination of the vertical velocity}\label{sec:vert-whoriz}

The profile of \cref{eq:vert-wlin} gives the vertical velocity at a single node; to obtain it at the
particle, the per-node profiles are combined horizontally over the one-ring of nodes surrounding the
particle. The combination method is selected by \conf{w\_interp}: \conf{nearest} (the profile at the
nearest node alone), \conf{leastsq} (a plane fit to the ring-node profiles, the default),
\conf{barycentric} (the vertex blend of \cref{ch:hvel} applied to the profiles), and \conf{idw}
(inverse-distance weighting). This least-squares stencil is over the surrounding \emph{nodes}, in
contrast to the horizontal velocity of \cref{sec:hvel-leastsq}, whose stencil is over the surrounding
\emph{element centres}; the two are distinct constructions that share a name. Only \conf{barycentric}
depends on the element type: on a quadrilateral it uses the Wachspress weights of
\cref{sec:hvel-quad}, on a triangle the triangular weights. The choice also decides whether a
three-dimensional run can use a fused kernel (\cref{ch:fused}).

\paragraph{z-level meshes.} When a profile is sampled at a neighbouring node, the evaluation depth is
capped at that neighbour's floor, $\depx_{\text{eval}}=\min(\depx,\,d_j)$, so that a shallow neighbour
cannot contribute a spurious below-floor velocity. The \conf{nearest} method evaluates the nearest
node's own profile without a cap: an RK stage position below that node's floor reads the profile below
the bottom, where the other methods' fallbacks cap the depth at the floor.

\paragraph{Sigma meshes.} On a sloping bottom a fixed depth cuts across sigma surfaces, and sampling
the neighbours at the particle's depth would mix the flow of different sigma levels. The neighbours
are therefore evaluated \emph{along the particle's sigma surface} (\cref{fig:vert-sigma}): with the
particle's fraction of its column
\begin{equation}\label{eq:vert-sigma-frac}
  \xi = \min\bigl(\max(\depx/H_p,\,0),\,1\bigr),
\end{equation}
neighbour $j$ is evaluated at $\depx_j=\xi\,H_j$, $H_j=\max(d_j+\eta_j,D_{\min})$. The particle column
$H_p$ is interpolated by the same method --- a least-squares plane (floored at $D_{\min}$), the
barycentric weights of the containing element, or inverse-distance weights --- which makes it exact for
a linearly sloping bottom under \conf{leastsq} and \conf{barycentric}; \conf{nearest} uses the nearest
node's column.

\begin{figure}[t]
\centering
\begin{tikzpicture}[scale=1.25]
  \fill[pattern=north east lines,pattern color=gray!45] (0,-2.0)--(5,-3.5)--(5,-3.8)--(0,-2.3)--cycle;
  \draw[laggreen,line width=0.9pt] (0,-2.0)--(5,-3.5);
  \draw[lagblue,line width=0.9pt] (0,0)--(5,0);
  \node[font=\footnotesize,lagblue,anchor=north west] at (0,-0.02) {surface};
  \node[font=\footnotesize,laggreen!70!black,anchor=north east] at (5,-3.55) {sloping floor};
  \draw[gray!60] (1.5,0)--(1.5,-2.45);
  \draw[gray!60] (4.0,0)--(4.0,-3.2);
  \draw[kelp!70,densely dashed,line width=0.8pt] (0,-1.2)--(5,-2.1);
  \node[font=\footnotesize,kelp!70!black,anchor=west] at (5.05,-2.1) {sigma surface $\xi=0.6$};
  \draw[gray,densely dotted,line width=0.8pt] (1.5,-1.47)--(5,-1.47);
  \node[font=\footnotesize,gray,anchor=west] at (5.05,-1.47) {depth $\depx$ (not used)};
  \node[drifter] at (1.5,-1.47) {};
  \node[meshnode] at (4.0,-1.92) {};
  \node[meshnodeopen] at (4.0,-1.47) {};
  \draw[<->,gray!80] (1.0,0)--(1.0,-2.3) node[pos=0.35,left,font=\footnotesize,text=black]{$H_p$};
  \draw[<->,gray!80] (4.3,0)--(4.3,-3.29) node[pos=0.8,right,font=\footnotesize,text=black]{$H_j$};
  \node[font=\footnotesize,anchor=north] at (4.0,-1.98) {$\xi H_j$};
  \node[font=\footnotesize,anchor=south] at (1.5,0.3) {particle column};
  \node[font=\footnotesize,anchor=south] at (4.0,0.3) {neighbour $j$};
\end{tikzpicture}
\caption{Sigma-following evaluation of the vertical velocity. The particle sits at the fraction
$\xi=\depx/H_p$ of its column; the neighbouring column is evaluated at $\xi H_j$ on the same sigma
surface (filled node), not at the particle's depth (open node), which on a sloping floor lies on a
different sigma level.}
\label{fig:vert-sigma}
\end{figure}

Every method falls back to the nearest-node profile when its stencil is unusable, so the combined
vertical velocity is never undefined.

\codeanchor{The horizontal combinations are \fn{interp\_w\_spatial} (z-level) and
\fn{interp\_w\_sigma\_spatial} (sigma), in \srcfile{numerics/vertical.py} and
\srcfile{numerics/vertical\_sigma.py}; the particle column is \fn{sigma\_column\_height} with the
per-method helpers \fn{\_H\_leastsq\_single}, \fn{\_H\_bary\_single} and \fn{\_H\_idw\_single}.}

\section{The step on a sigma mesh}\label{sec:vert-sigma-step}

FESOM-C's vertical velocity is a flux through the sigma interfaces, not the rate of the particle's
depth: with the surface rising and falling and the floor sloping, a particle at rest on its sigma
surface still changes depth. The step therefore integrates the particle's fraction of the water column,
\begin{equation}\label{eq:vert-sigma-ode}
  \xi = \frac{\depx}{H_p}, \qquad \frac{d\xi}{dt}=\frac{\Wvel}{H_p},
\end{equation}
with $\Wvel$ the negated file value of \cref{tab:vert-wsign} (including the direction factor). With
$\xi_0=\depx_n/H_p(\mathbf{x}_0,t_0)$ and $f_k=\Wvel_k/H_{p,k}$ evaluated at stage $k$,
\begin{align}
  \text{Euler:}\quad & \xi_{n+1}=\xi_0+\dt\,f_1, \label{eq:vert-sigma-euler}\\
  \text{RK2:}\quad & \xi_{1/2}=\xi_0+\tfrac{\dt}{2}f_1,\qquad \xi_{n+1}=\xi_0+\dt\,f_2, \label{eq:vert-sigma-rk2}\\
  \text{RK4:}\quad & \xi_{n+1}=\xi_0+\tfrac{\dt}{6}\,(f_1+2f_2+2f_3+f_4). \label{eq:vert-sigma-rk4}
\end{align}
Every stage places the particle at $\depx=\xi\,H_p(\mathbf{x}_k,t_k)$ in the column at the stage position
and time, and the end of the step at $\depx_{n+1}=\xi_{n+1}\,H_p(\mathbf{x}_{n+1},t_{n+1})$. The surface
elevation at the start and at the end of every step comes from one function of time --- the PCHIP
value where a four-snapshot stencil exists and PCHIP is enabled, otherwise the linear blend, and in
discrete mode the next record --- so the end of one step and the start of the next see the same surface
bit for bit, $\xi$ is continuous across steps, and a resumed run reproduces the uninterrupted one. With a
linear time interpolant, a linear bottom and a uniform dia-sigma velocity the discrete solution is
exact; with a dia-sigma velocity linear in $\xi$ the three schemes converge at orders one, two and four.
The layer of the horizontal velocity, the depth clamp and the diffusion kicks (\cref{ch:diffusion}) use
the same column $H_p$. The recorded vertical velocity on a sigma mesh is $\Wvel$ itself, the dia-sigma
flux, not the rate of the particle's depth.

\codeanchor{The sigma branches of \fn{\_euler\_step\_unfused}, \fn{\_rk2\_step\_unfused} and
\fn{\_rk4\_step\_unfused} with the helpers \fn{\_sigma\_col\_H} and \fn{\_sigma\_end\_eta}
(\srcfile{numerics/integrators.py}); the surface at the step boundaries is \fn{\_eta\_at\_time}
(\srcfile{tracking.py}). Sigma meshes always take the unfused path (\cref{ch:fused}).}

\section{Depth clamping at the surface and the floor}\label{sec:vert-clamp}

After the vertical position is advanced, the depth is clamped into the admissible range, and the
limits differ by family. On a z-level mesh the depth is clamped into $[\,z_0,\ d_{\text{nn}}\,]$ --- the
surface interface and the per-node sea floor of the nearest node --- and neither $\eta$ nor any
minimum-thickness parameter enters. On a sigma mesh the depth is clamped into
$[\,0,\ \max(H_p-z_{0b},\,0)\,]$, with $H_p$ the end-of-step column at the particle and $z_{0b}$ (default
\SI{0.1}{\metre}) holding the particle a fixed distance above the bed. The clamp is applied at every
step, so a particle that reaches the floor is not removed but tracks along it, in contrast to a particle
that leaves the horizontal domain, which is flagged and removed (\cref{ch:lifecycle}).

\codenote{The keys \conf{Dmin} ($D_{\min}$) and \conf{z0b\_min} ($z_{0b}$) are consumed only on sigma
meshes; the z-level clamp uses neither (nor $\eta$). $z_{0b}$ is a geometric clamp offset, not a
roughness length in any physics. The sigma surface clamp holds the particle at a literal $\depx=0$
rather than at $(1-\sigma_0)H$, so it coincides with the top interface only when $\sigma_0=1$ exactly.}

\codeanchor{Depth clamping is \fn{clamp\_depth} (z-level, limits $[z_0,d_{\text{nn}}]$) and
\fn{clamp\_depth\_sigma} (sigma, limits $[0,\max(H_p-z_{0b},0)]$) in the respective vertical modules,
invoked by \fn{\_clamp\_z} (\srcfile{numerics/integrators.py}); the defaults
$D_{\min}=\SI{0.5}{\metre}$, $z_{0b}=\SI{0.1}{\metre}$ are set in \srcfile{core/config.py}.}
